\documentclass[10pt,a4paper]{article}

\usepackage[utf8]{inputenc}
\usepackage[T1]{fontenc}
\usepackage[italian]{babel}
\usepackage[margin=2cm]{geometry}
\usepackage{microtype}
\usepackage{graphicx}
\usepackage{booktabs}
\usepackage{tabularx}
\usepackage{amsmath,amssymb}
\usepackage{xcolor}
\usepackage{enumitem}
\usepackage{float}
\usepackage[hidelinks]{hyperref}

\definecolor{accent}{HTML}{24566B}
\setlist{topsep=2pt, itemsep=2pt, parsep=0pt}
\setlength{\parindent}{0pt}
\setlength{\parskip}{4pt}

\newcommand{\file}[1]{\texttt{\small\detokenize{#1}}}
\newcommand{\shot}[2]{\fbox{\includegraphics[width=#1\textwidth,trim=240 85 315 58,clip]{img/#2.png}}}

\title{\color{accent}Studio esperti: proposta di task, materiali e prova}
\author{Proposta per il gruppo di ricerca --- NNModelling}
\date{19 novembre 2026}

\begin{document}
\maketitle

\section*{In breve}

\begin{itemize}
    \item \textbf{Tre task di difficoltà crescente}, svolti dagli esperti sia in
    NNModelling sia in PyTorch (Jupyter), in ordine bilanciato:
    (T1) lo stesso MLP degli studenti, (T2) self-attention a una testa,
    (T3) attention multi-testa con metavariabili e subflow.
    \item \textbf{Un solo foglio in inglese} per entrambi gli strumenti: una specifica
    neutrale per ogni task e due riquadri brevi (``In NNModelling'' / ``In PyTorch'')
    con il flusso di lavoro. Nessun addestramento in nessuno dei due strumenti.
    \item \textbf{Correzione automatica con gli stessi criteri} per i due strumenti:
    il progetto NNModelling viene eseguito in PyTorch dal runtime e sottoposto agli
    stessi controlli dei notebook.
    \item \textbf{Ho svolto io i tre task nell'app reale} (via VNC) partendo dai
    progetti \emph{start}: il correttore li ha promossi tutti e tre. Le osservazioni
    di usabilità emerse sono in Sezione~\ref{sec:usabilita}.
\end{itemize}

\section{Disegno dei task}

L'MLP identico a quello degli studenti consente il confronto diretto tra novizi ed
esperti; la crescita di difficoltà fornisce una curva di apprendimento; ogni task
parte da un proprio progetto, quindi un blocco non si propaga ai task successivi.
Il T3 introduce le funzioni avanzate scelte (metavariabili e subflow) e richiama
l'idea originale dell'LLM: la multi-head attention ne è il nucleo.

Tutti i task usano MNIST small: immagini \(x\in\mathbb{R}^{B\times1\times28\times28}\),
target \(y\in\{0,\dots,9\}^B\). Nei task con attention ogni immagine è una sequenza
di 28 righe (token) da 28 pixel, così si riusa lo stesso dataset e la stessa loss.

\begin{table}[H]
\centering\small
\begin{tabularx}{\textwidth}{@{}l>{\raggedright\arraybackslash}Xr@{}}
\toprule
Task & Architettura & Parametri \\
\midrule
T1 MLP & Flatten \(\to\) Linear(784,64) \(\to\) ReLU \(\to\) Linear(64,10); Cross Entropy & 50\,890 \\
T2 Attention, 1 testa & righe \([B,28,28]\) \(\to\) Linear(28,32)\(=H\);
\(Q,K,V=\)Linear(32,32); \(\mathrm{softmax}(QK^\top/\sqrt{32})V\); \(H+AV\);
Flatten \(\to\) Linear(896,10) & 13\,066 \\
T3 Multi-testa & come T2 con \texttt{d-model}\(=32\), \texttt{num-heads}\(=4\): ogni testa
\(Q_i,K_i,V_i=\)Linear\((d,d/h)\), scala \(1/\sqrt{d/h}\), teste concatenate &
\(3d^2+312d+10\) \\
\bottomrule
\end{tabularx}
\end{table}

\textbf{Scaffolding.} In T2 e T3 embedding, Add residuo e classificatore sono già
presenti in entrambi gli strumenti: l'esperto costruisce solo l'attention. Nel
progetto \emph{start} l'Add residuo è volutamente scollegato: un Add con un solo
input è valido e l'app mostrerebbe ``No model problems'' prima che il lavoro sia
iniziato.

\textbf{Stesse primitive.} In PyTorch sono vietati \texttt{nn.MultiheadAttention},
\texttt{F.scaled\_dot\_product\_attention} e \texttt{nn.Transformer*}: si usano solo
operazioni con un pacchetto equivalente in NNModelling (tabella sul foglio).
Senza causal mask (classificazione per righe, bidirezionale) restano sufficienti
i pacchetti \emph{core}.

\section{Materiali prodotti}

Tutto è nel repository, \textbf{non committato}:

\begin{table}[H]
\centering\small
\begin{tabularx}{\textwidth}{@{}lX@{}}
\toprule
Percorso & Contenuto \\
\midrule
\file{study/experts/tasks/expert-tasks.tex} & Foglio dei task (2 pagine, inglese) \\
\file{study/experts/projects/<task>/{start,solution}} & Progetti NNModelling di partenza e di riferimento \\
\file{study/experts/notebooks/} & Notebook dei partecipanti e soluzioni (\file{solutions/}) \\
\file{tools/expert_study/study_checks.py} & Controlli PyTorch, copiati accanto ai notebook (solo stdlib + torch) \\
\file{tools/expert_study/grade.py} & Correzione automatica di progetti e notebook (tabella, JSON, CSV) \\
\file{tools/expert_study/prepare_usb.py} & Prepara la pennetta: \file{master/}, \file{facilitator/}, \file{new_session.py} \\
\file{tools/expert_study/analyse_activity.py} & Separa il log di NNModelling per progetto (\file{split_activity_log.py}) e calcola i tempi per sotto-obiettivo con \file{tasks_expert_tN.json} \\
\file{tests/study_probe.c} & Analisi nativa a riga di comando (non salva mai il progetto) \\
\file{tests/test_expert_study_*.py} & 61 test automatici (\texttt{just test-expert-study}) \\
\file{docs/knowledge/testing/expert-study.md} & Specifica, criteri e risultati della prova \\
\bottomrule
\end{tabularx}
\end{table}

Le parti di codice (notebook, controlli, correttore, probe, pennetta) sono state
scritte da sottoagenti Sonnet e riviste da me. In review ho aggiunto un controllo:
la softmax deve normalizzare sull'ultima dimensione, altrimenti una softmax sulla
dimensione sbagliata passerebbe.

\section{Correzione automatica}

Gli stessi criteri valgono per il notebook e per il progetto NNModelling:

\begin{enumerate}
    \item forward su un batch fisso di 16 immagini: logits \([16,10]\) finiti; loss scalare
    uguale a \texttt{cross\_entropy(logits, y)};
    \item numero di parametri addestrabili pari alla tabella;
    \item nessun costrutto vietato (moduli e chiamate intercettate durante il forward);
    \item numero di mappe di attention \(28\times28\) per campione pari al numero di teste;
    \item solo T3: il modello viene ricostruito anche con \((d,h)=(64,8)\) e \((32,2)\):
    se un 32 è scritto a mano invece di usare la variabile, il controllo fallisce;
    \item solo NNModelling: analisi nativa completa e senza problemi; per il T3 esistono
    le metavariabili \texttt{d-model} e \texttt{num-heads}.
\end{enumerate}

Prova a vuoto della pennetta: una sessione nuova con i notebook risolti e i progetti
non toccati produce

{\small\begin{verbatim}
participant  tool     task  result  failed checks
P01          nnm      T1    FAIL    native analysis, logits, loss, parameters
P01          nnm      T2    FAIL    native analysis, logits, loss, parameters, attention heads
P01          nnm      T3    FAIL    d=32,h=4: native analysis, metavariables, ...
P01          pytorch  T1    PASS
P01          pytorch  T2    PASS
P01          pytorch  T3    PASS
\end{verbatim}}

mentre i tre progetti completati da me nella GUI danno \texttt{PASS} su T1, T2 e T3.

\section{Prova dei task nell'app}

Ho risolto ogni progetto \emph{start} nell'app Qt reale tramite VNC (click e
trascinamenti veri), seguendo solo il foglio. Le immagini a sinistra mostrano
ciò che vede l'esperto all'apertura; quelle a destra il progetto completato
(riordinato con \emph{Arrange} per la stampa).

\begin{figure}[H]
\centering
\shot{0.47}{t1-start}\hfill\shot{0.47}{t1-done}
\caption{T1 --- MLP. A sinistra lo \emph{start} (Input con MNIST small attivo, Output,
Loss Output); a destra il modello completo, senza problemi.}
\end{figure}

\begin{figure}[H]
\centering
\shot{0.47}{t2-start}\hfill\shot{0.47}{t2-done}
\caption{T2 --- self-attention a una testa. Lo \emph{start} lascia spazio tra
embedding e Add residuo; la scala è stata inserita con la formula del foglio
\texttt{\$(1/math.sqrt(32))}, risolta a \(0.1768\).}
\end{figure}

\begin{figure}[H]
\centering
\shot{0.47}{t3-start}\hfill\shot{0.47}{t3-done}
\caption{T3 --- multi-testa. Lo \emph{start} coincide con quello del T2. A destra la
radice completata: Horizontal Repeat con \texttt{times = \$num-heads}, embedding con
\texttt{\$d-model} e classificatore con \texttt{\$(28 * meta["d-model"])}.}
\end{figure}

\begin{figure}[H]
\centering
\shot{0.6}{t3-done-heads}
\caption{T3 --- interno dell'Horizontal Repeat: una testa con
\(Q,K,V\) = Linear(\texttt{\$d-model}, \texttt{\$(params.in\_features / meta["num-heads"])})
e scala \texttt{\$(1 / math.sqrt(meta["d-model"] / meta["num-heads"]))}, risolta a \(1/\sqrt8\).}
\end{figure}

\section{Osservazioni di usabilità emerse dalla prova}\label{sec:usabilita}

Escludo i limiti dello strumento VNC (niente gestore di finestre, socket rimasto
dopo l'arresto). Sono spunti da osservare negli esperti, non bug confermati:

\begin{itemize}
    \item \textbf{Nodi sovrapposti:} ogni nuovo nodo compare nello stesso punto, sopra gli
    altri, finché non si usa \emph{Arrange}.
    \item \textbf{Arrange e zoom:} con 14 nodi \emph{Arrange} rimpicciolisce molto schede e
    handle (difficili da agganciare) e ogni \emph{Arrange} reimposta lo zoom.
    \item \textbf{Pulsante che si sposta:} dopo aver selezionato un pacchetto, ``Add to graph''
    cambia posizione; il doppio click sul pacchetto lo aggiunge direttamente.
    \item \textbf{Inspector che scorre:} le righe si spostano quando un nodo ottiene un output
    valido o una formula aggiunge la riga ``Resolved''; il campo sotto il puntatore cambia.
    Anche \emph{Invio} su un parametro ha spostato la vista.
    \item \textbf{Trascinamenti mancati:} un trascinamento che manca l'handle sposta il nodo
    invece di creare l'arco.
    \item \textbf{Dialogo delle metavariabili:} il tipo predefinito è \emph{boolean} con valore
    \texttt{null}, e l'etichetta ``Value'' compare solo dopo aver scelto un altro tipo.
    \item \textbf{Formule (positivo):} \texttt{\$nome} e \texttt{\$(\ldots)} sono state accettate
    nello stesso campo, con il valore risolto mostrato subito sotto.
\end{itemize}

\section{Punti aperti per il gruppo}

\begin{enumerate}
    \item \textbf{SUS e NASA-TLX:} restano invariati. Se si compilano una sola volta alla
    fine, il giudizio su NNModelling arriva subito per il gruppo PyTorch\(\to\)NNM e dopo
    PyTorch per l'altro gruppo. Spostarli subito dopo la parte NNModelling non cambia i
    questionari; altrimenti va indicato tra i limiti.
    \item \textbf{Durata:} va misurata con un pilota. Stima: circa 5/10/10 minuti per task e
    per strumento per un esperto. Se il pilota supera l'ora, il T3 è il primo da accorciare.
    \item \textbf{Facilitatore:} tempo massimo per task, dopo il quale si passa al progetto
    successivo (le soluzioni sono in \file{facilitator/}), e regola sui suggerimenti.
    \item \textbf{Causal mask:} esclusa per restare nei pacchetti core; si può aggiungere al
    T2/T3 includendo il pacchetto dell'esempio LLM nei progetti \emph{start}.
    \item \textbf{Log di NNModelling:} i tre task vanno analizzati separatamente: gli \emph{start}
    di T2 e T3 sono identici, quindi il log va diviso per progetto prima di calcolare i tempi
    (\file{analyse_activity.py}); sulla mia prova tutti i 23 sotto-obiettivi risultano completati, ciascuno nel proprio task.
\end{enumerate}

\end{document}
